% !TEX root = ../main.tex
\begin{abstract}
% arXiv reproduces this abstract as plain text in the listing, so no custom
% macros and no \cite here, even though both would compile.
I argue that the assertion behaviour of large language models constitutes a
hyperintensional, non-normal modal operator. Knowledge and belief in the
received formal tradition are normal operators---standardly within the S4/S5
family and KD45 respectively---so language models are neither epistemic nor
doxastic systems in that sense. I define an output operator $\mathrm{O}_M$ and
a context-indexed grounding operator $\mathrm{G}_M$, and show the failure of
the rule RE and the axioms K, M, D, 4, and 5 for $\mathrm{O}_M$. The failure of
RE---congruentiality under logical equivalence---establishes
hyperintensionality. This rules out Kripke semantics for any accessibility
relation, and also rules out bare neighborhood semantics, which validates RE in
every frame; what remains is a hyperintensionalised neighborhood model, with
neighborhoods over formulations rather than over sets of worlds, or an
impossible-worlds semantics. The argument does not require that no
truth-correlated internal state exists, only that no such state governs
assertion.

Building on this apparatus I offer a second contribution: hallucination
redefined as ungrounded assertion, with no truth-value condition, and grounding
indexed to a practical context of jurisdiction, time, and intended use. The
resulting taxonomy of grounded against true shows that truth-matching
benchmarks misclassify both off-diagonal cells, systematically rewarding lucky
fabrication and penalising faithful transmission; I prove that accuracy places
no bound whatever on the grounding rate. Hallucination rate should instead be
measured as the conditional probability that an assertion is ungrounded given
that it is made, estimated by attribution audit---which is tractable only for
systems built to make their evidential base explicit at inference time.
Evaluation suites should include items whose answers are absent from the
accessible evidence, where the only correct behaviour is abstention, and should
penalise true answers on those items. I defend the resulting characterisation of
language models as sub-doxastic systems against interpretationist,
human-parallel, and deflationist objections, marking where the sub-doxastic
analogy gives out: there is no personal-level believer for these states to be
sub-personal relative to.
\end{abstract}
